\SourcePage{570}
\SGAsection{7}{附录：相交理论中的特殊化}
\begin{center}A. Grothendieck\end{center}
\begin{SGAresult}{\Term{lemma}{引理}}{7.1}
设 $i:Y\to X$ 是\Term{regular-immersion}{正则闭浸入}（第 VII 讲 1.4），其\Term{codimension}{余维数}\Term{bounded}{有界}，因而 $i$ 是\Term{perfect-morphism}{完美态射}（第 VII 讲 1.9），于是可以定义\Term{morphism}{态射}（第 IV 讲 2.12）
\[i^*:K_\bullet(X)\longrightarrow K_\bullet(Y).\]
另一方面，设 $N$ 为 $i$ 的\Term{conormal-sheaf}{余法层}；它是 $Y$ 上的\Term{finite-locally-free}{有限型局部自由层}（第 VII 讲 1.6），因而可以形成
\[\lambda_{-1}(N)\in K'(Y).\]
再考虑\Term{direct-image}{正像}\Term{homomorphism}{同态}（第 IV 讲 2.11）
\[i_*:K_\bullet(Y)\longrightarrow K_\bullet(X).\]
则有公式
\[i^*i_*(x)=x\lambda_{-1}(N)\qquad(x\in K_\bullet(Y)),\]
其中右边的乘积是相对于第 IV 讲 2.10 所定义的 $K_\bullet(Y)$ 的 $K'(Y)$ \Term{module-structure}{模结构}而言的。
\end{SGAresult}
\begin{proof}
证明与第 VII 讲 2.7 中在 $K'(X)$、$K'(Y)$ 内建立同一公式的证明实质上完全相同。
\end{proof}
\begin{SGAresult}{推论}{7.2}
设 $f:X\to S$ 是\Term{flat-morphism}{平坦态射}，$S$ 是\Term{trait}{离散赋值环谱}，即某个\Term{discrete-valuation-ring}{离散赋值环} $A$ 的谱；以 $s$ 表示 $S$ 的\Term{closed-point}{闭点}，$t$ 表示其\Term{generic-point}{泛点}，$i:X_s\to X$ 表示\Term{canonical}{典范}\Term{inclusion}{包含}。则 $i$ 的\Term{tor-dimension}{Tor 维数}\Term{finite}{有限}，因而可定义（第 IV 讲 2.12）
\[i^*:K_\bullet(X)\longrightarrow K_\bullet(X_s);\]
\SourcePage{571}
另一方面，有\Term{homomorphism}{同态}（第 IV 讲 2.11）
\[i_*:K_\bullet(X_s)\longrightarrow K_\bullet(X).\]
此时复合
\[i^*i_*:K_\bullet(X_s)\longrightarrow K_\bullet(X)\longrightarrow K_\bullet(X_s)\]
为零。
\end{SGAresult}
\begin{proof}
事实上，$i:X_s\to X$ 是\Term{codimension}{余维数}为 $1$ 的\Term{regular-immersion}{正则浸入}，其\Term{conormal-sheaf}{余法层} $N$ 是\Term{rank}{秩} $1$ 的\Term{free-module}{自由模}，因为 $X_s$ 可用一个整体方程 $f^*(g)=0$ 描述（其中 $g$ 是 $A$ 的\Term{uniformizer}{一致化参数}）。于是 $\lambda_{-1}(N)=0$，由 7.1 即得结论。
\end{proof}
以下始终假设 $X$ \Term{noetherian}{诺特}且在 $S$ 上\Term{flat-morphism}{平坦}。对 $X$ 及其\Term{open-subscheme}{开子概形} $X_t$ 和互补的\Term{closed-subscheme}{闭子概形} $X_s$ 应用\Term{homotopy-exact-sequence}{同伦正合列}（第 IX 讲 1.1\SGAsourcegap{原页在“IX 1.1”之后留有空白，未给出完整引注编号}）：
\[K_\bullet(X_s)\xrightarrow{i_*}K_\bullet(X)\xrightarrow{j^*}K_\bullet(X_t)\longrightarrow0,\]
其中 $j:X_t\to X$ 是\Term{inclusion}{包含}，亦即一个\Term{open-immersion}{开浸入}。利用 7.2 得到：
\begin{SGAresult}{命题}{7.3}
沿用 7.2 的记号，并假设 $X$ \Term{noetherian}{诺特}。存在唯一的\Term{group-homomorphism}{群同态}
\[\sigma_\bullet:K_\bullet(X_t)\longrightarrow K_\bullet(X_s)\tag{7.3.1}\]
（称为“\Term{specialization-homomorphism}{特殊化同态}”），使下图交换：
\[\begin{tikzcd}K_\bullet(X)\arrow[r,"j^*"]\arrow[d,"i^*"']&K_\bullet(X_t)\arrow[dl,"\sigma_\bullet"]\\K_\bullet(X_s)&\end{tikzcd}\tag{7.3.2}\]
\SourcePage{572}
借助\Term{ring-homomorphism}{环同态} $K'(X)\xrightarrow{j^*}K'(X_t)$ 和 $K'(X)\xrightarrow{i^*}K'(X_s)$（第 IV 讲 2.7 b），以及 $K_\bullet(X_t)$、$K_\bullet(X_s)$ 分别在 $K'(X_t)$、$K'(X_s)$ 上的\Term{module-structure}{模结构}（第 IV 讲 2.10），将 (7.3.1) 两端看作 $K'(X)$ \Term{module}{模}时，$\sigma_\bullet$ 是 $K'(X)$ \Term{module-homomorphism}{模同态}。
\end{SGAresult}
\begin{proof}
最后一个断言来自如下事实：由第 IV 讲 (2.12.2)，$i^*:K_\bullet(X)\to K_\bullet(X_s)$ 和\Term{epimorphism}{满同态} $j^*:K_\bullet(X)\to K_\bullet(X_t)$ 都是 $K'(X)$ \Term{module-homomorphism}{模同态}。
\end{proof}
下面逐一讨论\Term{specialization-homomorphism}{特殊化同态}的其他性质。
\textbf{7.4. 与\Term{filtered-category}{滤过的}相容性。}在各个 $K_\bullet$ 上赋予按\Term{support}{支集}\Term{dimension}{维数}定义的\Term{filtration}{滤过} $\operatorname{Fil}_\bullet K_\bullet$（1.1.1）。若 $f$ 为\Term{finite-type}{有限型}，更\Term{generic}{一般}地，若 $X$ 是\Term{catenary}{链状的}，则\Term{specialization-homomorphism}{特殊化同态}与这些\Term{filtration}{滤过}相容，因而诱导\Term{homomorphism}{同态}（也称为\Term{specialization-homomorphism}{特殊化同态}）
\[\sigma_\bullet:\operatorname{Gr}_\bullet(X_t)\longrightarrow\operatorname{Gr}_\bullet(X_s).\tag{7.4.1}\]
为此，由 1.14 只须证明：若 $Y_t$ 是 $X_t$ 的\Term{dimension}{维数}为 $d$ 的\Term{integral-scheme}{整子概形}，则 $\sigma_\bullet(\operatorname{cl}_\bullet(Y_t))$ 的\Term{filtration}{滤过}\Term{degree}{次数}不超过 $d$。设 $Y$ 为 $Y_t$ 在 $X$ 中的\Term{schematic-closure}{概形闭包}，于是 $Y$ 在 $S$ 上\Term{flat-morphism}{平坦}（EGA IV 2.8.1）。容易由此得到 $i^*(\operatorname{cl}_\bullet(Y))=\operatorname{cl}_\bullet(Y_s)$，两边的类分别取在 $K_\bullet(X)$ 和 $K_\bullet(X_s)$ 中，故
\[\sigma_\bullet(\operatorname{cl}_\bullet(Y_t))=\operatorname{cl}_\bullet(Y_s).\]
利用 $X$ 从而 $Y$ 为\Term{catenary}{链状的}这一事实，容易验证 $\dim Y_s\leq\dim Y_t$；当 $f$ 为\Term{finite-type}{有限型}时，将问题归结到 $X$ \Term{integral-scheme}{整}的情形后，这立刻由 EGA IV 7.1.13 得出。所述结论随即成立。
\SourcePage{573}
\textbf{7.5. “数的守恒”。}现在假设 $f$ \Term{proper-morphism}{固有}，于是 $X_t$、$X_s$ 分别在 $k(t)$、$k(s)$ 上\Term{proper-morphism}{固有}，故有 2.3 中的\Term{homomorphism}{同态}
\[\chi_t:K_\bullet(X_t)\to\mathbf Z,\qquad\chi_s:K_\bullet(X_s)\to\mathbf Z.\]
此时有关系
\[\chi_s(\sigma_\bullet(x))=\chi_t(x)\qquad(x\in K_\bullet(X_t))\tag{7.5.1}\]
（“数的守恒\Term{theorem}{定理}”），换言之，下图交换：
\[\begin{tikzcd}K_\bullet(X_t)\arrow[r,"\sigma_\bullet"]\arrow[dr,"\chi_t"']&K_\bullet(X_s)\arrow[d,"\chi_s"]\\&\mathbf Z\end{tikzcd}\tag{7.5.2}\]
为证明 (7.5.1)，仍可归结到 $x=\operatorname{cl}_\bullet(Y_t)$ 的情形，其中 $Y_t$ 是 $X_t$ 的\Term{integral-scheme}{整闭子概形}。沿用 7.4 证明中的记号，只须证明
\[\chi(Y_s)=\chi(Y_t),\]
而这是熟知的（EGA III 7.9.4，或第 III 讲 5.7）。
\textbf{7.6.} 显然有
\[\sigma_\bullet(\operatorname{cl}_\bullet(\mathcal O_{X_t}))=\operatorname{cl}_\bullet(\mathcal O_{X_s}).\tag{7.6.1}\]
结合 $\sigma_\bullet$ 的 $K'(X)$ \Term{linear}{线性}性，这说明下图交换：
\[\begin{tikzcd}&K'(X)\arrow[dl,"j^*"']\arrow[dr,"i^*"]&\\K'(X_t)\arrow[d]&&K'(X_s)\arrow[d]\\K_\bullet(X_t)\arrow[rr,"\sigma_\bullet"]&&K_\bullet(X_s)\end{tikzcd}\tag{7.6.2}\]
\SourcePage{574}
其中竖直\Term{arrow}{箭头}是\Term{canonical}{典范}\Term{homomorphism}{同态} $K'\to K_\bullet$（第 IV 讲 2.5.1）。
\textbf{7.7. $K'$ 与 $\operatorname{Gr}'$ 的\Term{specialization}{特殊化}。}当 $X_t$、$X_s$ \Term{regular-scheme}{正则}时（例如 $f$ \Term{smooth-morphism}{光滑}时），上述两个竖直\Term{arrow}{箭头}都是\Term{isomorphism}{同构}（第 IV 讲 2.5），于是\Term{specialization-homomorphism}{特殊化同态}\Term{solvable}{可解}释为\Term{homomorphism}{同态}
\[\sigma^\bullet:K'(X_t)\longrightarrow K'(X_s),\tag{7.7.1}\]
它也可描述为使下图交换的唯一\Term{homomorphism}{同态}：
\[\begin{tikzcd}K'(X)\arrow[r,"j^*"]\arrow[d,"i^*"']&K'(X_t)\arrow[dl,"\sigma^\bullet"]\\K'(X_s)&\end{tikzcd}\tag{7.7.2}\]
在此图中，$i^*$ 和 $j^*$ 是\Term{ring-homomorphism}{环同态}（第 IV 讲 2.7 b），且 $j^*$ \Term{surjective}{满射}，故\Term{specialization-homomorphism}{特殊化同态} (7.7.1) 是\Term{ring-homomorphism}{环同态}。同理，若 $X$ \Term{separated}{分离}，则\Term{homomorphism}{同态} (7.7.1) 甚至是 $\lambda$ \Term{ring-homomorphism}{环同态}。为说明这一点，只须注意：\Term{noetherian}{诺特}\Term{regular-scheme}{正则}\Term{separated}{分离}\Term{scheme}{概形}是\Term{divisorial}{除子的}（第 II 讲 2.2.7.1），因此它的 $K'$ 等同于其“朴素” $K'$（第 IV 讲 2.9 (ii)），后者自然带有 $\lambda$ \Term{ring-structure}{环结构}（第 VI 讲 3.2）。另一方面，$X_t$、$X_s$ \Term{regular-scheme}{正则}蕴含 $X$ 也\Term{regular-scheme}{正则}（EGA $0_{\mathrm{IV}}$ 17.1.8），而\Term{homomorphism}{同态} $i^*:K'(X)\to K'(X_s)$、$j^*:K'(X)\to K'(X_t)$ 是 $\lambda$ \Term{homomorphism}{同态}，故由 $i^*$ 取商所得的 (7.7.1) 亦然。
\SourcePage{575}
还须注意，$X$ \Term{regular-scheme}{正则}，因而\Term{normal-scheme}{正规}，所以其\Term{connected-component}{连通分支}与\Term{irreducible-component}{不可约分支}一致，而这些\Term{branch}{分支}都\Term{dominant-morphism}{支配} $S$；因此\Term{inclusion}{包含} $X_t\to X$ 诱导的\Term{homomorphism}{同态} $H^0(X,\mathbf Z)\to H^0(X_t,\mathbf Z)$ 是\Term{bijection}{双射}。于是，\Term{inclusion}{包含} $X_s\to X$ 诱导的\Term{homomorphism}{同态} $H^0(X,\mathbf Z)\to H^0(X_s,\mathbf Z)$ 也\Term{solvable}{可解}释为一个\Term{homomorphism}{同态}
\[H^0(X_t,\mathbf Z)\longrightarrow H^0(X_s,\mathbf Z),\tag{7.7.3}\]
它也应称为“\Term{specialization-homomorphism}{特殊化同态}”。此时直接由定义即可验证：考虑前述\Term{homomorphism}{同态} (7.7.3) 后，\Term{homomorphism}{同态} (7.7.1) 与两端在第 IV 讲 2.8 中定义的、分别以 $H^0(X_t,\mathbf Z)$、$H^0(X_s,\mathbf Z)$ 为底的\Term{augmented-algebra}{增广代数}结构相容。
仍假设 $X$ \Term{separated}{分离}，则上述结果特别说明 (7.7.1) 诱导一个\Term{graded-ring}{分次环}\Term{homomorphism}{同态}
\[\sigma^\bullet:\operatorname{Gr}^{\bullet}(X_t)\longrightarrow\operatorname{Gr}^{\bullet}(X_s),\tag{7.7.4}\]
并利用\Term{chern-class}{Chern 类}的概念（第 V 讲 6），得到\Term{commutative-diagram}{交换图}
\[\begin{tikzcd}K'(X_t)\arrow[r,"c_{X_t}"]\arrow[d]&\operatorname{Ch}\operatorname{Gr}^{\bullet}(X_t)\arrow[d]\\K'(X_s)\arrow[r,"c_{X_s}"]&\operatorname{Ch}\operatorname{Gr}^{\bullet}(X_s).\end{tikzcd}\]
\textbf{7.8. $\operatorname{Pic}$ 的\Term{specialization}{特殊化}。}保留 7.7 的假设：$X$ \Term{separated}{分离}，$X_t$、$X_s$ \Term{regular-scheme}{正则}。由\Term{canonical}{典范}\Term{isomorphism}{同构} (5.3.2)
\[\operatorname{Gr}^1\xrightarrow{\sim}\operatorname{Pic},\]
\Term{homomorphism}{同态} (7.7.4) 给出\Term{canonical}{典范}\Term{homomorphism}{同态}
\[\operatorname{Pic}(X_t)\longrightarrow\operatorname{Pic}(X_s).\tag{7.8.1}\]
\SourcePage{576}
还有一种\Term{equivalence}{等价}的描述：$X$ \Term{regular-scheme}{正则}意味着其\Term{local-ring}{局部环}是\Term{factorial-ring}{析因环}（Auslander--Buchsbaum \Term{theorem}{定理}，EGA IV 21.11.1），故 $X_t$ 上每个\Term{invertible-sheaf}{可逆层} $L_t$ 都可延拓为 $X$ 上的\Term{invertible-sheaf}{可逆层} $L$（EGA IV 21.6.11）。此时立即可验证，$\operatorname{cl}(L_t)\in\operatorname{Pic}(X_t)$ 在 (7.8.1) 下的像为 $\operatorname{cl}(L_s)\in\operatorname{Pic}(X_s)$，其中自然令 $L_s=L|_{X_s}$。
事实上，仅用 $X$ 的\Term{local-ring}{局部环}是\Term{factorial-ring}{析因环}以及 $X_s$ \Term{locally-integral}{局部整}这两个假设（无须\Term{regular-scheme}{正则}和\Term{separated}{分离}假设），就能证明 $\operatorname{cl}(L_s)$ 与所选的、将 $L_t$ 延拓为 $X$ 上\Term{invertible-sheaf}{可逆模}的 $L$ 无关，因此可直接得到\Term{homomorphism}{同态} (7.8.1)。为验证所述无关性，只须证明：若 $L_t\simeq\mathcal O_{X_t}$，则 $L_s\simeq\mathcal O_{X_s}$。这一假设意味着 $L\simeq\mathcal O_X(D)$，其中 $D$ 是 $X$ 上\Term{support}{支集}包含于 $X_s$ 的\Term{divisor}{除子}（参见 EGA IV 21.2.11）；立刻可归结到 $X_s$ \Term{integral-scheme}{整}的情形，于是 $D=nX_s$。但此时 $D=\operatorname{div}(g^n)$（$g$ 为 $A$ 的\Term{uniformizer}{一致化参数}），故 $\mathcal O_X(D)\simeq\mathcal O_X$，从而更有 $L_s\simeq\mathcal O_{X_s}$。
\textbf{7.9. \Term{specialization}{特殊化}与\Term{numerical-equivalence}{数值等价}。}设 $f$ \Term{proper-morphism}{固有}，$X_s$、$X_t$ \Term{regular-scheme}{正则}，并考虑两个\Term{specialization-homomorphism}{特殊化映射} (7.4.1)、(7.7.4)：
\[\begin{aligned}\sigma_\bullet&:\operatorname{Gr}_\bullet(X_t)\longrightarrow\operatorname{Gr}_\bullet(X_s),\\\sigma^\bullet&:\operatorname{Gr}^{\bullet}(X_t)\longrightarrow\operatorname{Gr}^{\bullet}(X_s).\end{aligned}\tag{7.9.1}\]
由于它们所来自的\Term{homomorphism}{同态}
\[\sigma:K'(X_t)\longrightarrow K'(X_s)\tag{7.9.2}\]
\SourcePage{577}
与\Term{ring-structure}{环结构}相容，再由 (7.5.1)，得到
\[\langle\sigma_\bullet(x_t),\sigma^\bullet(y_t)\rangle=\langle x_t,y_t\rangle\quad\bigl(x_t\in\operatorname{Gr}_\bullet(X_t),\ y_t\in\operatorname{Gr}^{\bullet}(X_t)\bigr),\tag{7.9.3}\]
其中采用 4.1 中 $\operatorname{Gr}_\bullet$ 与 $\operatorname{Gr}^{\bullet}$ 之间的\Term{pairing}{配对} $\langle x,y\rangle$。由此\Term{form}{形式}地\Term{pushout}{推出}：若 $x\in\operatorname{Gr}_\bullet(X_t)$（相应地，$x\in\operatorname{Gr}^{\bullet}(X_t)$）满足 $\sigma_\bullet(x)$（相应地，$\sigma^\bullet(x)$）\Term{numerical-equivalence}{数值等价}于零（4.5.1），则 $x$ 也如此。如果 (7.9.2) 在模去\Term{torsion}{挠元}后为\Term{epimorphism}{满同态}，则逆命题也成立；\Term{equivalence}{等价}地，取\Term{associated-graded}{相伴分次对象}所得的 (7.9.1) 中任一\Term{homomorphism}{同态}在模去\Term{torsion}{挠元}后为\Term{epimorphism}{满同态}。不过，这种情形仅在相当特殊时才出现。
\textbf{7.9.4.} 至少当 $f$ \Term{proper-morphism}{固有}且\Term{smooth-morphism}{光滑}时，可以相信上述逆命题总成立，即只要 $x$ \Term{numerical-equivalence}{数值等价}于零，$\sigma_\bullet(x)$（相应地，$\sigma^\bullet(x)$）也\Term{numerical-equivalence}{数值等价}于零。这将意味着\Term{homomorphism}{同态} (7.9.1) 在分别按\Term{numerical-equivalence}{数值等价}关系所得的\Term{quotient-group}{商群} $\operatorname{Gr}_\bullet^{\mathrm{num}}$、$\operatorname{Gr}^{\bullet}_{\mathrm{num}}$ 上诱导\Term{monomorphism}{单同态}。至少，上述结果说明 $\operatorname{Gr}_\bullet^{\mathrm{num}}(X_t)$ \Term{isomorphism}{同构}于 $\operatorname{Gr}_\bullet^{\mathrm{num}}(X_s)$ 的某个\Term{subgroup}{子群}的\Term{quotient-group}{商群}（特别地，前者的\Term{rank}{秩}至多等于后者的\Term{rank}{秩}）；对于 $\operatorname{Gr}^{\bullet}_{\mathrm{num}}$，以及对于 $K'_{\mathrm{num}}$ 本身，也有同样结论。
\textbf{7.10. 与\Term{direct-image}{正像}和\Term{inverse-image}{逆像}的相容性。}
设
\[g:X\longrightarrow Y\]
为 $S$ 上\Term{flat-morphism}{平坦}的\Term{noetherian}{诺特}\Term{scheme}{概形}之间的\Term{morphism}{态射}。以 $\sigma_X$（相应地，$\sigma_Y$）表示 $X$（相应地，$Y$）的\Term{specialization-homomorphism}{特殊化同态}，以 $i_X,j_X$（相应地，$i_Y,j_Y$）表示相应的 $i,j$，
\SourcePage{578}
即 $X$（相应地，$Y$）的特殊\Term{fiber}{纤维}和\Term{generic}{一般}\Term{fiber}{纤维}到整个\Term{scheme}{概形}中的\Term{morphism}{态射}。以
\[g_t:X_t\longrightarrow Y_t,\qquad g_s:X_s\longrightarrow Y_s\]
表示诱导的\Term{morphism}{态射}，于是有\Term{cartesian-square}{笛卡尔方块}
\[\begin{tikzcd}X_t\arrow[r,"g_t"]\arrow[d,"j_X"']&Y_t\arrow[d,"j_Y"]\\X\arrow[r,"g"]&Y\end{tikzcd}\qquad\begin{tikzcd}X_s\arrow[r,"g_s"]\arrow[d,"i_X"']&Y_s\arrow[d,"i_Y"]\\X\arrow[r,"g"]&Y.\end{tikzcd}\tag{7.10.0}\]
先设 $g$ \Term{proper-morphism}{固有}。我断言下图交换：
\[\begin{tikzcd}K_\bullet(X_t)\arrow[r,"g_{t*}"]\arrow[d,"\sigma_X"']&K_\bullet(Y_t)\arrow[d,"\sigma_Y"]\\K_\bullet(X_s)\arrow[r,"g_{s*}"]&K_\bullet(Y_s).\end{tikzcd}\tag{7.10.1}\]
这一结果推广了 7.5（因为当 $X=S$ 时，(7.3.2) 中各\Term{arrow}{箭头}均等同于 $\mathbf Z$ 的\Term{identity-map}{恒等映射}）。注意，若 $X$、$Y$ 满足 7.7 的条件，即\Term{fiber}{纤维}\Term{regular-scheme}{正则}，则 $g_t$、$g_s$ 自动具有\Term{finite}{有限}\Term{tor-dimension}{Tor 维数}，此图\Term{solvable}{可解}释为 $K'$ 的\Term{commutative-diagram}{交换图}；其改写留给读者。
为证明 (7.10.1) 的\Term{commutativity}{交换性}，利用 $\sigma_X$、$\sigma_Y$ 的定义，立刻可知只须证明下列棱柱中前面两个方块交换：
\[\begin{tikzcd}[row sep=large,column sep=large]K_\bullet(X_t)\arrow[rr,"g_{t*}"]\arrow[dd,dashed]&&K_\bullet(Y_t)\arrow[dd,dashed]\\&K_\bullet(X)\arrow[ul,"j_X^*"]\arrow[dl,"i_X^*"']\arrow[rr,"g_*"]&&K_\bullet(Y)\arrow[ul,"j_Y^*"']\arrow[dl,"i_Y^*"]\\K_\bullet(X_s)\arrow[rr,"g_{s*}"]&&K_\bullet(Y_s)\end{tikzcd}\tag{*}\]
\SourcePage{579}
而这两个方块正是由 (7.10.0) 的两个方块\Term{derive}{导出}的“\Term{base-change}{换底}”方块，因此可以应用第 IV 讲 3\SGAsourcegap{原页引注“IV 3”后未填具体编号}。只须验证它们对应于 $Y$ 上两对在 $Y$ 上\Term{tor-independent}{Tor 无关}的\Term{scheme}{概形} $(X,Y_t)$、$(X,Y_s)$。第一对的情形显然，因为 $Y_t\to Y$ \Term{flat-morphism}{平坦}；第二对也立即成立，因为\Term{immersion}{浸入} $Y_s\to Y$ 是\Term{codimension}{余维数} $1$ 的\Term{regular-immersion}{正则浸入}，且在\Term{base-change}{换底} $X\to Y$ 后仍然如此。
现在不再假设 $g$ \Term{proper-morphism}{固有}，而假设它有\Term{finite}{有限}\Term{tor-dimension}{Tor 维数}。我断言 $g_t$、$g_s$ 也如此。前者显然；后者来自刚才指出的 $X$ 和 $Y_s$ 在 $Y$ 上\Term{tor-independent}{Tor 无关}这一事实。于是可考虑\Term{homomorphism}{同态}方块
\[\begin{tikzcd}K_\bullet(Y_t)\arrow[r,"g_t^*"]\arrow[d,"\sigma_Y"']&K_\bullet(X_t)\arrow[d,"\sigma_X"]\\K_\bullet(Y_s)\arrow[r,"g_s^*"]&K_\bullet(X_s),\end{tikzcd}\tag{7.10.2}\]
其中水平\Term{arrow}{箭头}由第 IV 讲 2.12 定义。注意，若 $X$、$Y$ 满足 7.7 的条件，即\Term{fiber}{纤维}\Term{regular-scheme}{正则}，则（由于 $Y$ \Term{regular-scheme}{正则}）$g$ 自动具有\Term{finite}{有限}\Term{tor-dimension}{Tor 维数}，前图\Term{solvable}{可解}释为 $K'$ \Term{ring}{环}的一个\Term{commutative-diagram}{交换图}。
为证明 (7.10.2) 交换，仍只须证明与 (*) 类似的棱柱的前面方块交换，但把\Term{direct-image}{正像}\Term{homomorphism}{同态} $g_{t*},g_*,g_{s*}$ 换成\Term{inverse-image}{逆像}\Term{homomorphism}{同态} $g_t^*,g^*,g_s^*$。
\SourcePage{580}
考虑\Term{inverse-image}{逆像}的\Term{transitivity}{传递性}（第 IV 讲 2.12\SGAsourcegap{原页引注“IV 2.12”之后未填细分编号}），这些方块的\Term{commutativity}{交换性}立即由相应的方块 (7.10.0) 的\Term{commutativity}{交换性}得出。
\textbf{注记 7.10.3。}当 $g$ \Term{proper-morphism}{固有}时，(7.10.1) 的\Term{commutativity}{交换性}立即蕴含下图交换；它是取\Term{associated-graded}{相伴分次对象}所得的图：
\[\begin{tikzcd}\operatorname{Gr}_\bullet(X_t)\arrow[r,"g_{t*}"]\arrow[d,"\sigma_X"']&\operatorname{Gr}_\bullet(Y_t)\arrow[d,"\sigma_Y"]\\\operatorname{Gr}_\bullet(X_s)\arrow[r,"g_{s*}"]&\operatorname{Gr}_\bullet(Y_s).\end{tikzcd}\tag{7.10.1 bis}\]
同样，若 $X$、$Y$ 的\Term{fiber}{纤维}\Term{regular-scheme}{正则}且\Term{separated}{分离}，则 (7.10.2) 的\Term{commutativity}{交换性}蕴含下图交换：
\[\begin{tikzcd}\operatorname{Gr}^{\bullet}(Y_t)\arrow[r,"g_t^*"]\arrow[d,"\sigma_Y"']&\operatorname{Gr}^{\bullet}(X_t)\arrow[d,"\sigma_X"]\\\operatorname{Gr}^{\bullet}(Y_s)\arrow[r,"g_s^*"]&\operatorname{Gr}^{\bullet}(X_s).\end{tikzcd}\tag{7.10.2 bis}\]
又当 $X$、$Y$ 在 $S$ 上\Term{quasi-projective}{拟射影}且\Term{smooth-morphism}{光滑}时，已知 $g_t^*$、$g_s^*$ 与按\Term{support}{支集}\Term{codimension}{余维数}定义的\Term{filtration}{滤过}相容，于是得到类似 (7.10.2 bis) 的\Term{commutative-diagram}{交换图}，但将 $\operatorname{Gr}'$ 换为 $\operatorname{Gr}^{\bullet}_{\mathrm{top}}$（参见第 0 讲 2.1）。还可注意，若 $X$、$Y$ 在 $S$ 上\Term{quasi-projective}{拟射影}，且\Term{fiber}{纤维}\Term{regular-scheme}{正则}并\Term{separated}{分离}，则 $g_t$、$g_s$ 自动是\Term{complete-intersection-morphism}{完全交态射}（第 VIII 讲 1.1）；由相容性结果（第 VIII 讲 3.2），得到类似 (7.10.1 bis) 的\Term{commutative-diagram}{交换图}，但其中的 $\operatorname{Gr}_\bullet$ 换为 $\operatorname{Gr}^{\bullet}_{\mathbf Q}$。
\SourcePage{581}
\textbf{7.11.} 对\Term{algebraically-closed-field}{代数闭域} $k$ 上的\Term{scheme}{概形} $T$，可将 $K'(T)$ 中\Term{algebraic-equivalence}{代数等价}于零的元素组成的\Term{subgroup}{子群}定义为所有形如
\[E(x_1)-E(x_0)\tag{7.11.1}\]
的元素组成的\Term{subgroup}{子群}，其中 $x_0,x_1$ 是 $k$ 上一条\Term{proper-morphism}{固有}、\Term{regular-scheme}{正则}且\Term{connected}{连通}的\Term{algebraic-curve}{代数曲线} $C$ 上的两个 $k$ \Term{rational-point}{有理点}，$E\in K'(T\times_k C)$；对 $C$ 的任一 $k$ \Term{rational-point}{有理点} $x$，令
\[E(x)=i_x^*(E)\in K'(T),\qquad i_x:T\longrightarrow T\times_k C,\]
其中 $i_x$ 为由 $x$ 定义的\Term{morphism}{态射} $t\mapsto(t,x)$。形如 (7.11.1) 的元素确实构成 $K'(T)$ 的\Term{subgroup}{子群}，这可由一个标准论证得知：采用乘积 $C\times_k C'$ 及其点对 $(x_1,x'_1)$、$(x_0,x'_0)$，并利用如下事实：在 $k$ 上的\Term{quasi-projective}{拟射影}\Term{irreducible}{不可约}\Term{algebra}{代数}\Term{scheme}{概形}（这里是 $C\times_k C'$）中，总能作一条经过任意给定两点的\Term{irreducible}{不可约}\Term{algebraic-curve}{代数曲线}。（不愿接受此事实的读者，也可用形如 (7.11.1) 的元素生成的\Term{group}{群}来定义\Term{algebraic-equivalence}{代数等价}。）立即可知，\Term{algebraic-equivalence}{代数等价}于零的元素构成 $K'(T)$ 的一个\Term{ideal}{理想}；当 $\lambda^i$ 运算有定义时，即使用“朴素” $K'$ 时，该\Term{ideal}{理想}还在这些运算下\Term{stable}{稳定}。此外，显然该\Term{ideal}{理想}包含于 $K'(T)$ 的\Term{augmentation-ideal}{增广理想}。因此可定义\Term{quotient-ring}{商环}
\[K'_{\mathrm{alg}}(T),\tag{7.11.2}\]
使用朴素 $K'$ 时，它甚至是一个 $\lambda$ \Term{ring}{环}。
同样可定义 $K_\bullet(T)$ 中\Term{algebraic-equivalence}{代数等价}于零的元素\Term{subgroup}{子群}：在上面的定义中取 $E\in K_\bullet(T\times_k C)$，
\SourcePage{582}
并注意 $i_x:T\to T\times_k C$ 是\Term{perfect-morphism}{完美态射}，因为它是\Term{morphism}{态射} $j_x:x\to C$ 经\Term{flat-morphism}{平坦}\Term{base-change}{换底} $T\times_k C\to C$ 所得的，而由 $C$ \Term{regular-scheme}{正则}，$j_x$ 是\Term{perfect-morphism}{完美态射}。这样得到 $K_\bullet(T)$ 的\Term{quotient-group}{商群}
\[K_\bullet^{\mathrm{alg}}(T).\tag{7.11.3}\]
它自然是 $K'_{\mathrm{alg}}(T)$ \Term{module}{模}，因为利用 $i_x$ 的第 IV 讲 (2.12.2)，从定义立刻看出：$K'(T)$ 的一个元素与 $K_\bullet(T)$ 的一个元素的乘积，只要其中一个因子\Term{algebraic-equivalence}{代数等价}于零，就也\Term{algebraic-equivalence}{代数等价}于零。
（使用朴素 $K'$ 时）将 $\operatorname{Gr}^{\bullet}(T)$ 中\Term{algebraic-equivalence}{代数等价}于零的元素组成的\Term{ideal}{理想}定义为所有 $c^i(x)$（$i\geq1$，$x\in K'(T)$ \Term{algebraic-equivalence}{代数等价}于零）生成的\Term{ideal}{理想}；于是商\Term{graded-ring}{分次环}
\[\operatorname{Gr}^{\bullet}_{\mathrm{alg}}(T)\tag{7.11.4}\]
正是 $\lambda$ \Term{ring}{环} $K'_{\mathrm{alg}}(T)$ 的\Term{chern-ring}{Chern 环}。
最后，将 $\operatorname{Gr}_\bullet(T)$ 中\Term{algebraic-equivalence}{代数等价}于零的元素组成的\Term{graded-subgroup}{分次子群}（假设 $T$ \Term{locally-noetherian}{局部诺特}）定义为：其\Term{dimension}{维数} $i$ \Term{set-component}{分量}是 $\operatorname{Fil}_iK_\bullet(T)$ 中在 $K_\bullet(T)$ 内\Term{algebraic-equivalence}{代数等价}于零的元素\Term{subgroup}{子群}在 $\operatorname{Gr}_i(T)$ 中的像。立即可见，商\Term{graded-group}{分次群}
\[\operatorname{Gr}_\bullet^{\mathrm{alg}}(T)\tag{7.11.5}\]
由 $\operatorname{Gr}_\bullet(T)$ 在 $\operatorname{Gr}^{\bullet}(T)$ 上的\Term{module-structure}{模结构}（3.1）取商，成为\Term{ring}{环} $\operatorname{Gr}^{\bullet}_{\mathrm{alg}}(T)$ 上的\Term{module}{模}。
\textbf{7.11.6.} 不假设 $k$ \Term{algebraically-closed}{代数闭}，而在上述定义中要求 $C$ 在 $k$ 上\Term{geometrically-connected}{几何连通}，这些定义仍有意义。不过，这样的定义是否有用并不清楚；特别地，在 $\operatorname{Pic}(T)$ 的情形，似乎不能由此重新得到 2.4.1 所考虑的经典“几何”\Term{algebraic-equivalence}{代数等价}概念。
\SourcePage{583}
因此我们宁可规定：若 $K'(T)$ 的一个元素在 $K'(T\otimes_k\bar k)$ 中的像\Term{algebraic-equivalence}{代数等价}于零，则称它\Term{algebraic-equivalence}{代数等价}于零，其中 $\bar k$ 是 $k$ 的一个\Term{algebraic-closure}{代数闭包}；对\Term{group}{群} $K_\bullet(T)$、$\operatorname{Gr}^{\bullet}(T)$、$\operatorname{Gr}_\bullet(T)$ 中的\Term{algebraic-equivalence}{代数等价}也采用同一规定。
\textbf{7.12.} 有了这些定义，我断言：\Term{specialization-homomorphism}{特殊化同态} (7.3.1) 将\Term{algebraic-equivalence}{代数等价}于零的元素变为\Term{algebraic-equivalence}{代数等价}于零的元素，从而取商后定义\Term{homomorphism}{同态}
\[K_\bullet^{\mathrm{alg}}(X_t)\longrightarrow K_\bullet^{\mathrm{alg}}(X_s).\tag{7.12.1}\]
为此，利用 7.15 中说明的一个容易的相容性，立即可归结到 $A$ \Term{complete}{完备}且 $k(s)$ \Term{algebraically-closed}{代数闭}的情形。只须证明 $\sigma$ 将 $K_\bullet(X_t)$ 中形如 (7.11.1) 的元素变为 $K_\bullet(X_s)$ 中\Term{algebraic-equivalence}{代数等价}于零的元素。这里 $T=X_t$，因此将前面的 $C,E$ 写成 $C_t,E_t$。已知 $C_t$ 在 $k(t)$ 上\Term{projective}{射影}，故存在 $S$ 上\Term{projective}{射影}且\Term{flat-morphism}{平坦}的\Term{scheme}{概形}，其\Term{generic}{一般}\Term{fiber}{纤维}为 $C_t$（EGA IV 2.8.5）。利用 Abhyankar 关于算术\Term{surface}{曲面}的\Term{resolution-of-singularities}{奇点消解} [6]，甚至可选取 $C$ 为\Term{regular-scheme}{正则}\Term{scheme}{概形}。（注意，由 EGA IV 14.3.8，\Term{resolution-of-singularities}{消解}保持\Term{flat-morphism}{平坦}性。）由于 $C$ 在 $S$ 上\Term{proper-morphism}{固有}，$C_t$ 上的 $k(t)$ \Term{rational-point}{有理点} $x_{0t},x_{1t}$ 延拓为 $C$ 在 $S$ 上的\Term{section}{截面} $x_0,x_1$；且 $C$ \Term{regular-scheme}{正则}，故熟知 $C$ 在 $x_0(s),x_1(s)$ 处在 $S$ 上\Term{smooth-morphism}{光滑}（EGA IV 17.12.1 c)$\Rightarrow$a) 及 19.1.1）。
\SourcePage{584}
因此\Term{immersion}{浸入} $x_i:S\to X$（$i=0,1$）是\Term{regular-immersion}{正则浸入}，且经\Term{base-change}{换底} $s\to S$ 后仍为\Term{regular-immersion}{正则浸入} $x'_i:s\to C_s$。注意，由\Term{homotopy-exact-sequence}{同伦正合列}（第 IX 讲 1.1\SGAsourcegap{原页“IX 1.1”后未填完整引注编号}），$E_t\in K_\bullet(X_t\times_{k(t)}C_t)$ 来自某个 $E\in K_\bullet(X\times_S C)$，因而 $K_\bullet(X_t)$ 的元素 $E_t(x_{1t})-E_t(x_{0t})$ 来自 $K_\bullet(X)$ 的元素 $E(x_1)-E(x_0)$；由 (7.3.1)，它在 $K_\bullet(X_s)$ 中的像正是
\[i^*(E(x_1)-E(x_0)).\tag{*}\]
当然，对 $C$ 在 $S$ 上的\Term{section}{截面} $x$，如前令
\[E(x)=i_x^*(E),\qquad i_x:X\longrightarrow X\times_S C.\]
$i_x$ 可看作 $x$ 经\Term{flat-morphism}{平坦}\Term{base-change}{换底} $X\times_S C\to C$ 所得的\Term{morphism}{态射}，因此是\Term{regular-immersion}{正则浸入}，更是\Term{perfect-morphism}{完美态射}，故表达式 $E(x)$ 有意义。同理，表达式 $E_s(x_s)$ 也有定义；在下列由\Term{perfect-morphism}{完美态射}构成的\Term{commutative-diagram}{交换图}中应用\Term{inverse-image}{逆像}的\Term{transitivity}{传递性}（第 IV 讲 2.12）：
\[\begin{tikzcd}X\arrow[r,"i_x"]&X\times_S C\\X_s\arrow[u,"i"]\arrow[r,"i_{x_s}"']&X_s\times_{k(s)}C_s\arrow[u],\end{tikzcd}\]
得到 (*) 左边等于
\[E_s(x_{1s})-E_s(x_{0s}).\tag{**}\]
由于 $k(s)$ \Term{algebraically-closed}{代数闭}，$C_s$ 的\Term{reduced}{约化}\Term{irreducible-component}{不可约分支}在 $k(s)$ 上\Term{geometrically-integral}{几何整}，其交点也在 $k(s)$ 上有理。再利用 $C_s$ 的\Term{connected}{连通}性，容易得出 (**) 在 $K_\bullet(X_s)$ 中\Term{algebraic-equivalence}{代数等价}于零；
\SourcePage{585}
这里 $C_s$ 的\Term{connected}{连通}性来自 $C_t$ 的\Term{geometrically-connected}{几何连通}性。
由\Term{homomorphism}{同态} (7.12.1) 及 7.4，得到\Term{homomorphism}{同态}
\[\operatorname{Gr}_\bullet^{\mathrm{alg}}(X_t)\longrightarrow\operatorname{Gr}_\bullet^{\mathrm{alg}}(X_s).\tag{7.12.2}\]
若还假设 $X_s$、$X_t$ \Term{regular-scheme}{正则}，则 (7.12.1) 也\Term{solvable}{可解}释为\Term{ring-homomorphism}{环同态}
\[K'_{\mathrm{alg}}(X_t)\longrightarrow K'_{\mathrm{alg}}(X_s);\tag{7.12.3}\]
若 $X$ 又\Term{separated}{分离}，它是 $\lambda$ \Term{ring-homomorphism}{环同态}，从而定义\Term{homomorphism}{同态}
\[\operatorname{Gr}^{\bullet}_{\mathrm{alg}}(X_t)\longrightarrow\operatorname{Gr}^{\bullet}_{\mathrm{alg}}(X_s).\tag{7.12.4}\]
\textbf{7.12.5.} 设 $T$ 是域 $k$ 上\Term{smooth-morphism}{光滑}\Term{proper-morphism}{固有}的\Term{algebra}{代数}\Term{scheme}{概形}。容易验证，若 $x\in K'(T)$ 满足存在整数 $n\geq1$ 使 $nx$ \Term{algebraic-equivalence}{代数等价}于零，则 $x$ \Term{numerical-equivalence}{数值等价}于零。可以相信，至少 $T$ \Term{projective}{射影}时，逆命题总成立。在这一问题中，显然可假设 $k$ \Term{algebraically-closed}{代数闭}，并可将 $K'(T)$ 换成 $\operatorname{Gr}^{\bullet}(T)$ 或 $\operatorname{Gr}_\bullet(T)$，而不改变猜想。使用\Term{chow-ring}{Chow 环} $A(T)$（参见第 XIV 讲 4.1）也得到一个\Term{equivalence}{等价}的猜想；由此重新得到如下一个多少算是经典的猜想：若 $z$ 是 $T$ 上\Term{numerical-equivalence}{数值等价}于零的\Term{algebra}{代数}\Term{cycle}{圈}，则存在整数 $n\geq1$ 使 $nz$ \Term{algebraic-equivalence}{代数等价}于零。若此猜想对 $X_t$ 正确，则由\Term{homomorphism}{同态} (7.3.1) 与\Term{algebraic-equivalence}{代数等价}的相容性结果，\Term{pushout}{推出}\Term{specialization-homomorphism}{特殊化同态}（当 $f:X\to S$ 为\Term{smooth-morphism}{光滑}\Term{proper-morphism}{固有态射}时）确实与\Term{numerical-equivalence}{数值等价}相容，正如 7.9 所猜测的那样。
\SourcePage{586}
\textbf{7.13. 与 $\ell$ 进\Term{cohomology}{上同调}的关系。}
回顾（SGA 5 第 VII 讲）：对\Term{scheme}{概形} $T$ 上任一\Term{finite-locally-free}{有限型局部自由模} $E$，以及不同于 $T$ 的各\Term{residue-characteristic}{剩余特征}的任一\Term{prime-number}{素数} $\ell$，定义有 $\ell$ 进\Term{chern-class}{Chern 类}（$i\geq0$）
\[c^i(E)\in H^{2i}(T,\mathbf Z_\ell(i)),\]
它们满足可\Term{additive}{加性}公式
\[c(E)=c(E')c(E'')\]
当 $E$ 是\Term{locally-free}{局部自由模} $E'$ 被\Term{locally-free}{局部自由模} $E''$ 所作的\Term{extension}{扩张}时成立。由此可\Term{additive}{加性}，\Term{total-chern-class}{总 Chern 类} $c(E)=\sum_i c^i(E)$ 来自一个\Term{homomorphism}{同态}
\[c:K'(T)\longrightarrow\widehat G(H^{2*}(T,\mathbf Z_\ell(*)))=1+\prod_{i\geq1}H^{2i}(T,\mathbf Z_\ell(i)),\tag{7.13.0}\]
其中右边是一个\Term{multiplicative-group}{乘法群}，如第 V 讲 1 所述，$K'(T)$ 是“朴素” $K'$。由此得到\Term{canonical-map}{典范映射}
\[c^i:K'(T)\longrightarrow H^{2i}(T,\mathbf Z_\ell(i)).\tag{7.13.1}\]
设 $X$ 是 $S$ 上\Term{proper-morphism}{固有}且\Term{flat-morphism}{平坦}的\Term{scheme}{概形}，各\Term{fiber}{纤维}\Term{regular-scheme}{正则}，并设 $S$ 是\Term{henselian}{Hensel 的}。假设 $\ell$ 不等于 $S$ 的\Term{residue-characteristic}{剩余特征}。于是可作图
\[\begin{tikzcd}K'(X_t)\arrow[r,"c^i"]\arrow[d,"\sigma"']&H^{2i}(X_t,\mathbf Z_\ell(i))\\K'(X_s)\arrow[r,"c^i"]&H^{2i}(X_s,\mathbf Z_\ell(i))\arrow[u,"\sigma'"']\end{tikzcd}\tag{7.13.2}\]
其中水平\Term{arrow}{箭头}是分别对应于 $X_t$、$X_s$ 的“\Term{chern-class}{Chern 类}”映射，而 $\sigma'$ 是复合\Term{homomorphism}{同态}
\[H(X_s)\xleftarrow[\sim]{i^*}H(X)\xrightarrow{j^*}H(X_t);\]
\SourcePage{587}
这里写出的\Term{arrow}{箭头}是\Term{restriction}{限制}\Term{homomorphism}{同态}；由\Term{proper-base-change}{固有换底定理}（SGA 4 第 XII 讲 5.5），$i^*$ 是\Term{isomorphism}{同构}。我断言方块 (7.13.2) 交换。证明与建立 (7.10.2) 的\Term{commutativity}{交换性}所用的证明实质上相同；这里只须对\Term{morphism}{态射} $j:X_t\to X$、$i:X_s\to X$，利用\Term{chern-class}{Chern 类}的形成与\Term{inverse-image}{逆像}相容这一事实。
\textbf{7.13.3.} 注意，方块 (7.13.2) 是类似方块的\Term{projective-limit}{投射极限}，其中系数为 $\ell$ 进\Term{sheaf}{层} $\mathbf Z_\ell(i)$ 的\Term{cohomology}{上同调}换成系数为\Term{torsion-sheaf}{挠层} $\mu_n^{\otimes i}$ 的\Term{cohomology}{上同调}（$n$ 遍取 $\ell$ 的幂）。因此，(7.13.2) 的\Term{commutativity}{交换性}实际上只是说：对 $\ell$ 的每个幂 $n$，涉及\Term{module}{模} $n$ 的\Term{chern-class}{Chern 类}的方块交换。选取 $K=k(t)$ 的一个\Term{algebraic-closure}{代数闭包} $\bar K$，由此得到 $k=k(s)$ 的一个\Term{algebraic-closure}{代数闭包} $\bar k$（即 $V$ 在 $\bar K$ 中的\Term{normal-closure}{正规闭包}的\Term{residue-field}{剩余域}）。对 $\bar K$ 中每个\Term{finite}{有限}子\Term{extension}{扩张} $K_\alpha$，可以考虑相似的\Term{commutative-diagram}{交换图}（固定 $n$），它对应于由 $f:X\to S$ 经\Term{base-change}{换底} $S_\alpha\to S$ 所得的情形，其中 $S_\alpha$ 是 $S$ 在 $K_\alpha$ 中的\Term{normalization}{正规化}。注意，由 Krull--Akizuki \Term{theorem}{定理}和 $S$ 的\Term{henselian}{Hensel 性}，$S_\alpha$ 确实是一个\Term{trait}{离散赋值环谱}；此外，这里必须假设 $X$ 在 $S$ 上的\Term{fiber}{纤维}\Term{geometrically-regular}{几何正则}，即 $X$ 在 $S$ 上\Term{smooth-morphism}{光滑}，以使该假设经\Term{base-change}{换底}后保持。这样，随 $\alpha$ 变化得到由\Term{commutative-square}{交换方块}构成的\Term{inductive-system}{归纳系}。取\Term{inductive-limit}{归纳极限}，并利用 SGA 4 第 VII 讲 5.8 和第 IV 讲 3.2\SGAsourcegap{原页“IV 3.2”后未给出完整引注编号}，得到\Term{commutative-square}{交换方块}（参见 7.16）：
\SourcePage{588}
\[\begin{tikzcd}K'(X_{\bar t})\arrow[r,"c^i"]\arrow[d,"\sigma"']&H^{2i}(X_{\bar t},\mu_n^{\otimes i})\\K'(X_{\bar s})\arrow[r,"c^i"]&H^{2i}(X_{\bar s},\mu_n^{\otimes i})\arrow[u,"\sigma'"'] .\end{tikzcd}\]
其中 $X_{\bar t}$、$X_{\bar s}$ 是 $X$ 的\Term{geometric-fiber}{几何纤维}。此外，现在已知\Term{cohomology}{上同调}的\Term{specialization-homomorphism}{特殊化同态} $\sigma'$ 是\Term{isomorphism}{同构}（SGA 4 第 XVI 讲 2.2），因此前一方块的\Term{commutativity}{交换性}也\Term{representable-functor}{可表示}为将该\Term{arrow}{箭头}倒转所得方块的\Term{commutativity}{交换性}。在所得的\Term{commutative-square}{交换方块}中再对 $n$ 取\Term{projective-limit}{投射极限}，得到\Term{commutative-square}{交换方块}
\[\begin{tikzcd}K'(X_{\bar t})\arrow[r]\arrow[d,"\sigma"']&H^{2i}(X_{\bar t},\mathbf Z_\ell(i))\arrow[d,"\sim"]\\K'(X_{\bar s})\arrow[r]&H^{2i}(X_{\bar s},\mathbf Z_\ell(i)).\end{tikzcd}\tag{7.13.4}\]
\textbf{7.13.5.} 这一\Term{commutativity}{交换性}也可表述如下：若 $x_t\in K'(X_t)$，$x_s$ 为其在 $K'(X_s)$ 中的像，则 $c^i(x_t)$、$c^i(x_s)$ 是 $R^{2i}f_*\mathbf Z_\ell(i)$ 的一个\Term{section}{截面}分别在\Term{geometric-point}{几何点} $\bar t$、$\bar s$ 处的值。此\Term{section}{截面}也很容易解释：取 $x\in K'(X)$，其\Term{restriction}{限制}为 $x_t$，则 $c^i(x)\in H^{2i}(X,\mathbf Z_\ell(i))$ 定义了 $R^{2i}f_*\mathbf Z_\ell(i)$ 的所求\Term{section}{截面}。
\textbf{7.13.6.} 设 $T$ 为\Term{scheme}{概形}。取\Term{associated-graded}{相伴分次对象}后，\Term{homomorphism}{同态} (7.13.0) 定义\Term{homomorphism}{同态}
\[\psi_i:\operatorname{Gr}^i(T)\longrightarrow H^{2i}(T,\mathbf Z_\ell(i)),\]
其中元素 $x\in\operatorname{Fil}^iK'(T)$ 的像 $\psi_i(x)$ 按定义为 $c^i(x)$。（(7.13.0) 与 $K'(T)$ 上的 $\lambda$ \Term{filtration}{滤过}以及右边的\Term{degree}{次数}\Term{filtration}{滤过}相容，
\SourcePage{589}
这容易由第 V 讲 6.6.4 及下述事实得出：将 $K'(T)$ 看作\Term{binomial-ring}{二项式环} $H^0(T,\mathbf Z)$ 上的\Term{augmented-algebra}{增广} $\lambda$ \Term{algebra}{代数}时，$c^i$ 定义了第 V 讲 6.13 意义下的 $K'(T)$ 的\Term{chern-theory}{Chern 理论}。）由此得到映射
\[\operatorname{cl}:\operatorname{Gr}^i(T)\longrightarrow H^{2i}(T,\mathbf Q_\ell(i))\overset{\mathrm{def}}=H^{2i}(T,\mathbf Z_\ell(i))\otimes_{\mathbf Z_\ell}\mathbf Q_\ell,\tag{7.13.7}\]
其公式为
\[\operatorname{cl}(x^{(i)})=\frac{1}{(-1)^{i-1}(i-1)!}\,\psi_i(x^{(i)}).\tag{7.13.8}\]
参见第 XIV 讲 5.1、5.2。于是，图 (7.13.2)、(7.13.4) 的\Term{commutativity}{交换性}显然蕴含类似图的\Term{commutativity}{交换性}：
\[\begin{tikzcd}\operatorname{Gr}^i(X_t)_{\mathbf Q}\arrow[r,"\operatorname{cl}"]\arrow[d,"\sigma_{\mathbf Q}"']&H^{2i}(X_t,\mathbf Q_\ell(i))\\\operatorname{Gr}^i(X_s)_{\mathbf Q}\arrow[r,"\operatorname{cl}"]&H^{2i}(X_s,\mathbf Q_\ell(i))\arrow[u,"\sigma'"'] ,\end{tikzcd}\tag{7.13.9}\]
\[\begin{tikzcd}\operatorname{Gr}^i(X_{\bar t})_{\mathbf Q}\arrow[r,"\operatorname{cl}"]\arrow[d]&H^{2i}(X_{\bar t},\mathbf Q_\ell(i))\arrow[d,"\sim"]\\\operatorname{Gr}^i(X_{\bar s})_{\mathbf Q}\arrow[r,"\operatorname{cl}"]&H^{2i}(X_{\bar s},\mathbf Q_\ell(i)).\end{tikzcd}\tag{7.13.10}\]
\textbf{注记 7.13.11。}最后两个图的\Term{commutativity}{交换性}可以表述为：\Term{algebra}{代数}\Term{cycle}{圈}的\Term{specialization-homomorphism}{特殊化}与相应 $\ell$ 进\Term{cohomology-class}{上同调类}的形成相容。不过，这种表述并不十分令人满意，因为我们不得不把系数 $\mathbf Z_\ell(i)$ 换成 $\mathbf Q_\ell(i)$，也就是说“模去了\Term{torsion}{挠元}”。为得到一个不忽略\Term{torsion}{挠元}的陈述，
\SourcePage{590}
似乎必须使用\Term{chow-ring}{Chow 环}，因为正是在该\Term{ring}{环}上，\Term{homomorphism}{同态} $\operatorname{cl}$ 才自然地定义为取值于 $H^{2*}(-,\mathbf Z_\ell(*))$ 的\Term{homomorphism}{同态}；参见第 XIV 讲 4.1。遗憾的是，目前尚未发展出\Term{chow-ring}{Chow 环}的\Term{specialization-homomorphism}{特殊化}理论（参见 7.14），使我们能用\Term{chow-ring}{Chow 环}以及以 $\mathbf Z_\ell(*)$ 为系数的\Term{cohomology}{上同调}写出类似 (7.13.9)、(7.13.10) 的\Term{commutative-square}{交换方块}。当 $f:X\to S$ \Term{smooth-morphism}{光滑}且\Term{proper-morphism}{固有}时，只能说有\Term{commutative-square}{交换方块}
\[\begin{tikzcd}Z^i(X_t)\arrow[r,"\operatorname{cl}"]\arrow[d,"\sigma"']&H^{2i}(X_t,\mathbf Z_\ell(i))\\Z^i(X_s)\arrow[r,"\operatorname{cl}"]&H^{2i}(X_s,\mathbf Z_\ell(i))\arrow[u] ,\end{tikzcd}\tag{7.13.12}\]
它类似于 (7.13.9)；还有一个类似 (7.13.10) 的相应方块，其写法留给读者。这里 $Z^i$ 表示\Term{codimension}{余维数}为 $i$ 的\Term{algebra}{代数}\Term{cycle}{圈}\Term{group}{群}（不模去任何关系），水平\Term{arrow}{箭头}是 SGA 5 第 IV 讲定义的“相伴\Term{cohomology-class}{上同调类}”\Term{homomorphism}{同态}；$\sigma$ 表示\Term{algebra}{代数}\Term{cycle}{圈}的\Term{specialization-homomorphism}{特殊化同态}，它将一个\Term{cycle}{圈} $Z_t$ 送到\Term{cycle}{圈} $Z_s$，其中 $Z$ 是 $Z_t$ 在 $X$ 中的\Term{schematic-closure}{概形闭包}（参见 7.15）。这一\Term{commutativity}{交换性}来自 SGA 5 第 IV 讲建立的如下事实：若 $Z$ 是 $X$ 的\Term{closed-subscheme}{闭子概形}，在每个\Term{fiber}{纤维}上的\Term{codimension}{余维数}为 $i$，则可赋予它一个\Term{cohomology-class}{上同调类} $\operatorname{cl}(Z)\in H^{2i}(X,\mathbf Z_\ell(i))$，其形成与任意\Term{base-change}{换底}交换；特别地，它在 $X_t$（相应地，$X_s$）上的\Term{restriction}{限制}等于 $\operatorname{cl}(Z_t)$（相应地，$\operatorname{cl}(Z_s)$）。当然，若已知 (7.13.12) 的\Term{homomorphism}{同态} $\sigma$ 将\Term{rational-equivalence}{有理等价}于零的\Term{cycle}{圈}变为\Term{rational-equivalence}{有理等价}于零的\Term{cycle}{圈}（$X$ 在 $S$ 上\Term{projective}{射影}），
\SourcePage{591}
则可在该图中将 $Z^i$ 换为 $A^i$，得到以 $X_t$、$X_s$ 的\Term{chow-ring}{Chow 环}表述的所需类型的图。
\textbf{注记 7.13.13。}设 $T$ 是\Term{algebraically-closed-field}{代数闭域}上的\Term{smooth-morphism}{光滑}\Term{projective}{射影}\Term{scheme}{概形}。由 SGA 5 第 IV 讲所发展的\Term{form}{形式}体系，显然有：若 $Z$ 是 $T$ 上\Term{codimension}{余维数}为 $i$ 的\Term{cycle}{圈}，且其在 $H^{2i}(T,\mathbf Q_\ell(i))$ 中的像为零，则 $Z$ \Term{numerical-equivalence}{数值等价}于零。逆命题也成立，是\Term{algebra}{代数}\Term{cycle}{圈}理论的标准猜想之一。由第 VII 讲 4.11，并使用定义 (7.13.6)，立即可知将\Term{algebra}{代数}\Term{cycle}{圈}\Term{group}{群}换成\Term{group}{群} $\operatorname{Gr}^{\bullet}(T)_{\mathbf Q}$ 得到一个\Term{equivalence}{等价}猜想；也可换成\Term{group}{群} $K'(T)$，表述为：$K'(T)$ 的一个元素\Term{numerical-equivalence}{数值等价}于零，当且仅当它的\Term{augmentation}{增广}及其在 $H^{2i}(T,\mathbf Q_\ell(i))$ 中的各\Term{chern-class}{Chern 类}为零。
于是，由 (7.13.10) 的\Term{commutativity}{交换性}立即得到：若 $X_{\bar t}$ 满足上述猜想，则\Term{specialization-homomorphism}{特殊化同态} $\operatorname{Gr}^{\bullet}(X_{\bar t})\to\operatorname{Gr}^{\bullet}(X_{\bar s})$（从而更有 $\operatorname{Gr}^{\bullet}(X_t)\to\operatorname{Gr}^{\bullet}(X_s)$）与\Term{numerical-equivalence}{数值等价}相容，如 7.9 所猜测的那样。还应注意，此处采用的猜想（\Term{numerical-equivalence}{数值等价}蕴含 $\mathbf Q_\ell$ \Term{cohomological-equivalence}{上同调等价}）弱于 7.11.5 所考虑的猜想。
\textbf{7.14. 圈群及\Term{cycle-class-group}{圈类群}的\Term{specialization-homomorphism}{特殊化同态}。}
若 $T$ 是\Term{noetherian}{诺特}\Term{scheme}{概形}，以 $Z_i(T)$（相应地，$Z^i(T)$）表示 $T$ 上纯\Term{dimension}{维数} $i$（相应地，纯\Term{codimension}{余维数} $i$）的\Term{algebra}{代数}\Term{cycle}{圈}\Term{group}{群}。当 $f:X\to S$ 为\Term{finite-type}{有限型}时，定义\Term{specialization-homomorphism}{特殊化同态}
\SourcePage{592}
\[\sigma_i:Z_i(X_t)\longrightarrow Z_i(X_s)\tag{7.14.1}\]
如下：对 $X_t$ 的每个纯\Term{dimension}{维数}为 $i$ 的素\Term{cycle}{圈} $Z_t$，将其等同于 $X_t$ 的一个\Term{integral-scheme}{整子概形}，定义 $\sigma_i(Z_t)\in Z_i(X_s)$ 为与 $X_s$ 的\Term{closed-subscheme}{闭子概形} $Z_s$ 相伴的\Term{cycle}{圈}；其中 $Z$ 是 $Z_t$ 在 $X$ 中的\Term{schematic-closure}{概形闭包}（它是 $X$ 的一个在 $S$ 上\Term{flat-morphism}{平坦}的\Term{closed-subscheme}{闭子概形}，见 EGA IV 2.8），其\Term{generic}{一般}\Term{fiber}{纤维}为给定的 $Z_t$，并满足 $\dim Z_s=\dim Z_t$（EGA IV 7.1.13）。当 $X_s$ \Term{locally-equidimensional}{局部等维}（即过任一点的所有\Term{irreducible-component}{不可约分支}\Term{dimension}{维数}相同）时，立即可验证刚定义的\Term{specialization-homomorphism}{特殊化同态}诱导\Term{homomorphism}{同态}
\[\sigma^i:Z^i(X_t)\longrightarrow Z^i(X_s).\tag{7.14.1 bis}\]
现在假设 $f$ \Term{quasi-projective}{拟射影}，从而 $X_t$、$X_s$ 上的\Term{cycle}{圈}有\Term{rational-equivalence}{有理等价}的概念（参见 [5]，那里是在一个\Term{algebraically-closed-field}{代数闭域}上建立该理论的）。很有理由相信，\Term{homomorphism}{同态} (7.14.1) 将\Term{rational-equivalence}{有理等价}于零的\Term{cycle}{圈}变为\Term{rational-equivalence}{有理等价}于零的\Term{cycle}{圈}，因而在\Term{cycle-class}{圈类}的“Chow”\Term{group}{群}上诱导相应的\Term{homomorphism}{同态}
\[\sigma_i:A_i(X_t)\longrightarrow A_i(X_s),\tag{7.14.2}\]
或者，当 $X_s$ \Term{locally-equidimensional}{局部等维}时，诱导
\[\sigma^i:A^i(X_t)\longrightarrow A^i(X_s).\tag{7.14.2 bis}\]
可以考虑用仿照 7.12 的方法证明这一点，不过此处不必再用\Term{resolution-of-singularities}{奇点消解}，也不必归结到\Term{residue-field}{剩余域}\Term{algebraically-closed}{代数闭}的情形，
\SourcePage{593}
因为这里有 $C=\mathbf P^1_S$ 可作为使\Term{cycle}{圈}变动的“\Term{parameter-curve}{参数曲线}”。然而会遇到困难：沿用 7.12 的记号，$E$ 现在应表示 $\mathbf P^1_S\times_S X$ 上的一个\Term{cycle}{圈}，使得 $E_t(x_{1t})-E_t(x_{0t})$ “有定义”，但这丝毫不意味着\Term{cycle}{圈} $E_s(x_{1s})-E_s(x_{0s})$ 也有定义。因此正如预料，为定义 (7.14.2)，仍须提出一个 Chow 的“\Term{moving-lemma}{移动引理}”型的命题。
不过应注意，(7.4.1) 的存在表明 (7.14.1) 与\Term{surjective}{满射}\Term{homomorphism}{同态} $Z_\bullet(T)\to\operatorname{Gr}_\bullet(T)$ 在\Term{cycle}{圈}上定义的\Term{equivalence-relation}{等价关系}相容，而这一关系仅比\Term{linear-equivalence}{线性等价}稍粗；事实上，若 $T$ 是域上的\Term{quasi-projective}{拟射影}\Term{smooth-morphism}{光滑}\Term{scheme}{概形}，我们在别处指出，两者仅相差\Term{torsion}{挠元}（第 XIV 讲 4.1、4.2）。因此在这种情形，得到类似 (7.14.2)，或者更确切地说类似 (7.14.2 bis) 的\Term{homomorphism}{同态}，不过两端都须与 $\mathbf Q$ 作\Term{tensor-product}{张量积}。
要直接定义\Term{homomorphism}{同态} (7.14.2)，最自然的方式似乎是仿照 7.3。但这要求有一个适用于比域上的\Term{quasi-projective}{拟射影}\Term{scheme}{概形}更\Term{generic}{一般}的\Term{scheme}{概形}、并涵盖 $X$ 的 Chow \Term{group}{群} $A_i$ 理论（在\Term{regular-scheme}{正则}情形则是\Term{chow-ring}{Chow 环}理论），以及这些\Term{group}{群}的一个\Term{contravariant}{反变}规则，使我们能定义到 $X_s$ 的\Term{restriction}{限制}\Term{homomorphism}{同态} $A_i(X)\to A_i(X_s)$。参见第 XIV 讲 8。一旦有了这样充分\Term{generic}{一般}的 Chow \Term{group}{群}及\Term{chow-ring}{Chow 环}理论，就应能将前面列举的关于 $K$ 理论中的\Term{specialization-homomorphism}{特殊化同态}的全部结果改写到 Chow 理论中。
\SourcePage{594}
\textbf{7.15. 改变 $S$。}设 $S'$ 是另一个\Term{trait}{离散赋值环谱}，
\[S'\longrightarrow S\]
是\Term{trait}{离散赋值环谱}之间的\Term{morphism}{态射}，即一个\Term{local-morphism}{局部}\Term{dominant-morphism}{支配态射}（即把\Term{closed-point}{闭点}映成\Term{closed-point}{闭点}，把\Term{generic-point}{泛点}映成\Term{generic-point}{泛点}）。以 $s'$（相应地，$t'$）表示 $S'$ 的\Term{closed-point}{闭点}（相应地，\Term{generic-point}{泛点}），令 $X'=X\times_S S'$。假设 $X'$ \Term{noetherian}{诺特}（例如 $X$ 在 $S$ 上为\Term{finite-type}{有限型}）。有\Term{commutative-diagram}{交换图}
\[\begin{tikzcd}X_s\arrow[r,"i_X"]&X&X_t\arrow[l,"j_X"']\\X'_{s'}\arrow[u]\arrow[r,"i_{X'}"']&X'\arrow[u]&X'_{t'}\arrow[u]\arrow[l,"j_{X'}"]\end{tikzcd}\]
其中连接上下两行的\Term{arrow}{箭头}是\Term{flat-morphism}{平坦态射}。因此可沿图中所有\Term{arrow}{箭头}取 $K_\bullet$ 的\Term{inverse-image}{逆像}\Term{homomorphism}{同态}（第 IV 讲 2.12）。应用定义 7.3，得到\Term{commutative-diagram}{交换图}
\[\begin{tikzcd}K_\bullet(X_t)\arrow[r,"\sigma_X"]\arrow[d]&K_\bullet(X_s)\arrow[d]\\K_\bullet(X'_{t'})\arrow[r,"\sigma_{X'}"]&K_\bullet(X'_{s'}),\end{tikzcd}\tag{7.15.1}\]
其中竖直\Term{arrow}{箭头}是\Term{inverse-image}{逆像}\Term{homomorphism}{同态}。立即可由此得到将 $K_\bullet$ 换成 $\operatorname{Gr}_\bullet$ 的类似\Term{commutative-diagram}{交换图}；当 $X_s$、$X_t$ \Term{regular-scheme}{正则}时，还可将 $K_\bullet$、$\operatorname{Gr}_\bullet$ 分别换成 $K'$、$\operatorname{Gr}'$。
\SourcePage{595}
\textbf{7.16. 转到\Term{geometric-fibre}{几何纤维}。}设 $\bar K$ 为 $K=k(t)$ 的一个\Term{algebraic-closure}{代数闭包}，它是其\Term{finite}{有限}子\Term{extension}{扩张} $K_i$ 的\Term{inductive-limit}{归纳极限}；设 $\bar A=\varinjlim\bar A_i$ 是 $A$ 在 $\bar K$ 中的\Term{normalization}{正规化}，它是 $A$ 在 $K_i$ 中的\Term{normalization}{正规化} $\bar A_i$ 的\Term{inductive-limit}{归纳极限}。选取 $\bar A$ 的一个\Term{maximal-ideal}{极大理想} $\bar{\mathfrak m}$，它在各个 $\bar A_i$ 上诱导\Term{maximal-ideal}{极大理想} $\mathfrak m_i$，令 $B_i=(\bar A_i)_{\mathfrak m_i}$。由 Krull--Akizuki \Term{theorem}{定理}，这仍是\Term{discrete-valuation-ring}{离散赋值环}；以 $S_i$ 表示其谱，$s_i,t_i$ 表示其各点。\Term{residue-field}{剩余域} $\bar k=\bar A/\bar{\mathfrak m}$ 是 $k=k(s)$ 的一个\Term{algebraic-closure}{代数闭包}，并且是 $B_i$ 的各\Term{residue-field}{剩余域} $k_i=\bar A_i/\mathfrak m_i=k(s_i)$ 的\Term{inductive-limit}{归纳极限}。令 $\bar t=\operatorname{Spec}(\bar K)$、$\bar s=\operatorname{Spec}(\bar k)$。现在假设 $X$ 在 $S$ 上为\Term{finite-type}{有限型}，于是 $X_{\bar t}$、$X_{\bar s}$ 分别在 $\bar t$、$\bar s$ 上为\Term{finite-type}{有限型}，因而\Term{noetherian}{诺特}。由第 IV 讲 3.2.3，有
\[K_\bullet(X_{\bar t})\xleftarrow{\sim}\varinjlim_i K_\bullet(X_{t_i}),\qquad K_\bullet(X_{\bar s})\xleftarrow{\sim}\varinjlim_i K_\bullet(X_{s_i}).\]
另一方面，由 7.15，有从各 $K_\bullet(X_{t_i})$ 构成的\Term{inductive-system}{归纳系}到各 $K_\bullet(X_{s_i})$ 构成的\Term{inductive-system}{归纳系}的一个\Term{homomorphism}{同态}。取\Term{limit}{极限}后，得到\Term{canonical}{典范}\Term{homomorphism}{同态}
\[K_\bullet(X_{\bar t})\longrightarrow K_\bullet(X_{\bar s}),\tag{7.16.1}\]
它自然仍称为\Term{specialization-homomorphism}{特殊化同态}。由 7.4，并利用\Term{functor}{函子} $\operatorname{Gr}_\bullet$（同 $K_\bullet$ 一样）与此处的\Term{limit}{极限}交换这一事实（其验证留给读者），它诱导\Term{homomorphism}{同态} $\operatorname{Gr}_\bullet(X_{\bar t})\to\operatorname{Gr}_\bullet(X_{\bar s})$。同样，若 $f$ \Term{smooth-morphism}{光滑}，则可将 (7.16.1) 解释为 $K'$ 的\Term{homomorphism}{同态}；它还是\Term{ring-homomorphism}{环同态}，并且当 $X$ \Term{separated}{分离}时甚至是 $\lambda$ \Term{ring-homomorphism}{环同态}，此时它诱导 $\operatorname{Gr}'$ 的\Term{homomorphism}{同态}。
\SourcePage{596}
\textbf{7.16.2.} 前面各节对“算术”\Term{fiber}{纤维}证明的\Term{specialization-homomorphism}{特殊化同态}的全部性质，在\Term{geometric-fiber}{几何纤维}情形中的推广，这里不再逐一讨论；它们均可由取\Term{limit}{极限}直接得到。仅指出 7.9.4 的一个有趣结论：当 $f$ \Term{smooth-morphism}{光滑}且\Term{proper-morphism}{固有}时，$\operatorname{Gr}_\bullet^{\mathrm{num}}(X_{\bar t})$ \Term{isomorphism}{同构}于 $\operatorname{Gr}_\bullet^{\mathrm{num}}(X_{\bar s})$ 的某个\Term{subgroup}{子群}的\Term{quotient-group}{商群}；特别地，其\Term{rank}{秩}不超过后一个\Term{group}{群}的\Term{rank}{秩}。因此，这是一个关于\Term{module}{模}\Term{numerical-equivalence}{数值等价}的\Term{cycle}{圈}\Term{group}{群}的\Term{rank}{秩}的\Term{semicontinuity}{半连续性}结果。
\textbf{7.17. 任意底的情形。}现在设
\[f:X\longrightarrow S\]
为\Term{noetherian}{诺特}\Term{prescheme}{预概形}之间的\Term{flat-morphism}{平坦}\Term{quasi-compact-morphism}{拟紧态射}。设 $s,t$ 为 $S$ 的两点，$s$ 是 $t$ 的\Term{specialization}{特殊化}；我们希望联系两个\Term{fiber}{纤维} $X_t,X_s$ 以及相应\Term{geometric-fiber}{几何纤维}的 $K_\bullet$ 等\Term{group}{群}。为此注意（EGA II 7.1.9），存在\Term{morphism}{态射}
\[g:S'\longrightarrow S,\]
其中 $S'$ 是\Term{trait}{离散赋值环谱}，$g$ 将\Term{closed-point}{闭点} $s'$ 映成 $s$，将\Term{generic-point}{泛点} $t'$ 映成 $t$；甚至可假设 $k(t)\xrightarrow{\sim}k(t')$。假设 $X'=X\times_S S'$ \Term{noetherian}{诺特}；当 $f$ 为\Term{finite-type}{有限型}，或者 $S'$ 在 $S$ 上\Term{essentially-finite-type}{本质有限型}时，这一条件成立（EGA IV 1.3.8）；若 $S$ 本身为\Term{universally-japanese}{普遍日本的}，就可以作后一假设（EGA $0_{\mathrm{IV}}$ 23.1.1）。于是 $X'$ 的\Term{specialization-homomorphism}{特殊化同态}定义\Term{homomorphism}{同态}
\[\sigma:K_\bullet(X_t)\xrightarrow{\sim}K_\bullet(X_{t'})\longrightarrow K_\bullet(X_{s'}).\tag{7.17.1}\]
\SourcePage{597}
这样得到 $K_\bullet(X_t)$ 上的一个\Term{specialization-homomorphism}{特殊化同态}，但它取值于 $K_\bullet(X_{s'})$，而不是 $K_\bullet(X_s)$；其中 $X_{s'}$ 由 $X_s$ 经\Term{base-field}{基域}的\Term{extension}{扩张} $k(s)\to k(s')$ 得到（若 $S$ 为\Term{universally-japanese}{普遍日本的}，可以假设该\Term{extension}{扩张}为\Term{finite-type}{有限型}）。如 7.16 所述转到\Term{geometric-fiber}{几何纤维}，同样得到\Term{homomorphism}{同态}
\[\sigma:K_\bullet(X_{\bar t})\longrightarrow K_\bullet(X_{\bar s'}).\tag{7.17.2}\]
所以这里“几何”\Term{specialization-homomorphism}{特殊化同态}的值仍不是落在 $K_\bullet(X_{\bar s})$ 中，而是落在 $K_\bullet(X_{\bar s'})$ 中，其中 $X_{\bar s'}$ 由 $X_{\bar s}$ 经\Term{algebraically-closed-field}{代数闭域}的\Term{extension}{扩张} $k(\bar s)\to k(\bar s')$ 得到。
\textbf{7.17.3.} \Term{generic}{一般}而言，\Term{homomorphism}{同态} $K_\bullet(X_{\bar s})\to K_\bullet(X_{\bar s'})$ 不是\Term{isomorphism}{同构}（因为\Term{group}{群} $K_\bullet(X_{\bar s})$ 中可能存在“\Term{continuous}{连续}”部分）。不过，容易验证，模去\Term{algebraic-equivalence}{代数等价}后它是\Term{isomorphism}{同构}；若 $f$ \Term{smooth-morphism}{光滑}且\Term{proper-morphism}{固有}，则模去\Term{numerical-equivalence}{数值等价}（4.5.1）后也如此。因此利用 7.12，得到 (7.17.2) 诱导\Term{specialization-homomorphism}{特殊化同态}
\[K_\bullet^{\mathrm{alg}}(X_{\bar t})\longrightarrow K_\bullet^{\mathrm{alg}}(X_{\bar s});\tag{7.17.3.1}\]
若 $f$ \Term{proper-morphism}{固有}且\Term{smooth-morphism}{光滑}，并以\Term{specialization-homomorphism}{特殊化同态} (7.17.2) 与\Term{numerical-equivalence}{数值等价}相容为前提（参见 7.9.4），则同样得到\Term{homomorphism}{同态}
\[K_\bullet^{\mathrm{num}}(X_{\bar t})\longrightarrow K_\bullet^{\mathrm{num}}(X_{\bar s}).\tag{7.17.3.2}\]
相应的 $\operatorname{Gr}_\bullet$、$K'$、$\operatorname{Gr}'$ 构造，以及从前面各节发展的性质直接得到的这些\Term{specialization-homomorphism}{特殊化同态}的性质，留给读者在需要时展开。这里只指出，
\SourcePage{598}
由 7.16.2 可知，$\operatorname{Gr}_\bullet^{\mathrm{num}}(X_{\bar t})$ \Term{isomorphism}{同构}于 $\operatorname{Gr}_\bullet^{\mathrm{num}}(X_{\bar s})$ 的某个\Term{subgroup}{子群}的\Term{quotient-group}{商群}（$f$ \Term{proper-morphism}{固有}且\Term{smooth-morphism}{光滑}）；特别地，其\Term{rank}{秩}仍不超过后一个\Term{group}{群}的\Term{rank}{秩}。
\textbf{7.17.4.} 当 $S$ 在 $s$ 处\Term{regular-scheme}{正则}且 $t$ 为\Term{maximal-point}{极大点}时，可在 $S$ 于 $s$ 处的\Term{localization}{局部化}中以 $s$ 为\Term{center}{中心}作\Term{blowup}{爆破}，再在\Term{blowup}{爆破}中 $s$ 的\Term{fiber}{纤维}的\Term{generic-point}{泛点}处\Term{localization}{局部化}，来构造 $S'$。这样得到的 $k(s')$ 是 $k(s)$ 的一个\Term{pure-transcendental-extension}{纯超越扩张}。不难证明：若 $k'$ 是域 $k$ 的\Term{pure-transcendental-extension}{纯超越扩张}，则对 $k$ 上任一\Term{noetherian}{诺特}\Term{scheme}{概形} $T$，\Term{canonical}{典范}\Term{homomorphism}{同态} $K_\bullet(T)\to K_\bullet(T\otimes_k k')$ 是\Term{isomorphism}{同构}；证明时将 $k'$ 表为形如 $\operatorname{Spec}k[T_1,\ldots,T_n]$ 的\Term{scheme}{概形}的\Term{affine-open}{仿射开集}的\Term{affine}{仿射}\Term{ring}{环}的\Term{inductive-limit}{归纳极限}，并利用第 IX 讲 1.6\SGAsourcegap{原页“IX 1.6”后未填完整引注编号}。因此由此得出，(7.17.1) 可视为\Term{canonical}{典范}\Term{homomorphism}{同态}
\[K_\bullet(X_t)\longrightarrow K_\bullet(X_s).\tag{7.17.4.1}\]
但须注意，不能由此\Term{derive}{导出}\Term{homomorphism}{同态} $K_\bullet(X_{\bar t})\to K_\bullet(X_{\bar s})$，因为转到 $K=k(t)$ 的一个\Term{finite-extension}{有限扩张} $K_i$ 并将 $S$ 在 $K_i$ 中\Term{normalization}{正规化}后，$S$ 的\Term{regular-scheme}{正则}性假设将不再保持。
另一种定义\Term{homomorphism}{同态} (7.17.4.1) 的方式，是先归结到 $S$ 为以 $s$ 为\Term{closed-point}{闭点}的\Term{local-scheme}{局部概形}的情形（这很容易），再考虑 $S$ 的一列递增的\Term{regular-scheme}{正则}\Term{closed-subscheme}{闭子概形} $S_i$，其\Term{generic-point}{泛点}为 $t_i$，使 $\dim S_i=i$（$0\leq i\leq n=\dim S$）。对每个 $0\leq i\leq n-1$，$\mathcal O_{S_{i+1},t_i}$ 是\Term{residue-field}{剩余域}为 $k(t_i)$、\Term{fraction-field}{分式域}为 $k(t_{i+1})$ 的\Term{discrete-valuation-ring}{离散赋值环}，所以有\Term{specialization-homomorphism}{特殊化同态}
\[K_\bullet(X_{t_{i+1}})\longrightarrow K_\bullet(X_{t_i});\]
\SourcePage{599}
这些\Term{homomorphism}{同态}对 $0\leq i\leq n-1$ 的复合定义\Term{homomorphism}{同态} (7.17.4.1)，而且可以像 7.3 那样证明，它与前面用\Term{blowup}{爆破}定义的\Term{homomorphism}{同态}相同。不过，第二种方法的优点在于，这里甚至有形如 (7.16.1) 的\Term{homomorphism}{同态}：
\[K_\bullet(X_{\bar t_{i+1}})\longrightarrow K_\bullet(X_{\bar t_i}),\]
其中对各 $k(t_i)$ 的\Term{algebraic-closure}{代数闭包} $k(\bar t_i)$ 作适当相容的选取。将它们复合，即可定义\Term{homomorphism}{同态}
\[K_\bullet(X_{\bar t})\longrightarrow K_\bullet(X_{\bar s}),\tag{7.17.4.2}\]
它无疑本质上依赖于 $S_i$ 的选择，也就是说依赖于 $t_i$ 的选择。
\textbf{7.17.5.} 即使不再假设 $S$ 在 $s$ 处\Term{regular-scheme}{正则}且 $t$ 为\Term{maximal-point}{极大点}，只要对于 $S$ 中以 $t$ 为\Term{generic-point}{泛点}的\Term{integral-scheme}{整}\Term{closed-subscheme}{闭子概形} $T$，或仅对其在 $s$ 处的\Term{localization}{局部化}，有\Term{resolution-of-singularities}{奇点消解}，就能推广\Term{specialization-homomorphism}{特殊化同态} (7.17.4.2) 的构造。由 Hironaka 的结果，特别当 $\mathcal O_{S,s}$ 为\Term{characteristic-zero}{特征零}的\Term{excellent-ring}{优良环}时，此假设成立。事实上，此时可以找到\Term{morphism}{态射} $S'\to\operatorname{Spec}(\mathcal O_{T,s})$，其中 $S'$ \Term{regular-scheme}{正则}；取\Term{fiber}{纤维} $S'_s$ 的一个\Term{closed-point}{闭点}，即可对 $S'$ 上的 $X'=X\times_S S'$、$S'$ 的\Term{generic-point}{泛点} $t'$ 及其\Term{closed-point}{闭点} $s'$ 应用前述结果。由此得到形如 (7.17.4.2) 的\Term{homomorphism}{同态}，因为 $k(s')$ 是 $k(s)$ 的\Term{finite-extension}{有限扩张}，因而具有相同的\Term{algebraic-closure}{代数闭包}。
\textbf{7.17.6.} 在第 7.17 节考虑的所有构造中，为构造\Term{specialization-homomorphism}{特殊化同态}，都归结到了底为\Term{trait}{离散赋值环谱}的情形。
\SourcePage{600}
因此，前面各节在底为\Term{trait}{离散赋值环谱}时逐一讨论的\Term{specialization-homomorphism}{特殊化同态}的性质，自动推广到底更\Term{generic}{一般}的情形。当然，所述定义方法同样适用于 $\operatorname{Gr}_\bullet$、$K'$（当 $f$ \Term{smooth-morphism}{光滑}时）及 $\operatorname{Gr}'$（当 $f$ 还\Term{separated}{分离}时）。
\textbf{7.17.7.} 最后指出，\Term{homomorphism}{同态} (7.17.3.1)、(7.17.3.2) 或\Term{homomorphism}{同态} (7.17.4.2) 的定义，本质上依赖于某些任意选择：在 7.17.3 中，选择 $S$ 之上的适当\Term{trait}{离散赋值环谱}；在 7.17.4 中，选择一列 $S_i$；在 7.17.5 中，还要选择一个\Term{resolution-of-singularities}{奇点消解}。$\operatorname{Gr}_\bullet$ \Term{equivariant}{等变}式亦然。很容易构造例子：令 $s$ 为一条\Term{algebraic-curve}{代数曲线}的\Term{ordinary-double-point}{普通二重点}，$X$ 为 $S$ 上\Term{projective}{射影}\Term{smooth-morphism}{光滑}的\Term{scheme}{概形}（例如相对\Term{dimension}{维数}为 $2$ 的\Term{abelian-scheme}{阿贝尔概形}），则 $\operatorname{Gr}^1=\operatorname{Pic}$ 甚至 $\operatorname{Gr}^1_{\mathrm{alg}}=\operatorname{NS}$（\Term{neron-severi-group}{Néron--Severi 群}）上的\Term{specialization-homomorphism}{特殊化同态}确实依赖于在 $S$ 之上所选的\Term{trait}{离散赋值环谱} $S'$；这里\Term{equivalence}{等价}地说，它依赖于在 $S$ 的\Term{normalization}{正规化}中选择 $s$ 上方的一个点。
\textbf{7.18. 附录的附录：某些\Term{abelian-group}{交换群}值\Term{functor}{函子}的\Term{base-field-extension}{基域扩张}。}
这里详述 7.17.4 中提出的一个技术问题。
\textbf{7.18.1.} 设 $k$ 为域，$C$ 为\Term{integral-algebra}{整 $k$ 代数}及这些\Term{algebra}{代数}间的\Term{flat-homomorphism}{平坦同态}构成的\Term{category}{范畴}，$F$ 为 $C$ 上取值于\Term{set}{集合}\Term{category}{范畴} $(\mathrm{Ens})$ 的\Term{functor}{函子}。假设 $F$ 与\Term{filtered-inductive-limit}{滤过归纳极限}交换。设 $K$ 为 $k$ 的\Term{extension}{扩张}，我们寻求使 $F(k)\to F(K)$ 为\Term{injective}{单射}（相应地，\Term{surjective}{满射}）的判别准则。
\SourcePage{601}
\textbf{7.18.2.} 假设处于以下两种情形之一：
\begin{enumerate}
\item[a)] 域 $k$ \Term{separably-closed}{可分闭}，且 $K/k$ 是\Term{separable-extension}{可分扩张}。
\item[b)] 域 $k$ 无限，且 $K$ 是其\Term{pure-transcendental-extension}{纯超越扩张}。
\end{enumerate}
则映射 $F(k)\to K(K)$ 为\Term{injective}{单射}。
事实上，利用 $K$ 是 $k$ 上各\Term{finite-type}{有限型}\Term{subalgebra}{子代数}的\Term{filtered-inductive-limit}{滤过归纳极限}、其过渡\Term{morphism}{态射}为\Term{flat-morphism}{平坦}这一事实，以及 $F$ 与\Term{filtered-inductive-limit}{滤过归纳极限}交换，在 a) 中只须证明：若 $A$ 是 $k$ 上的\Term{finite-type}{有限型}\Term{integral-algebra}{整代数}，且其\Term{fraction-field}{分式域} $K$ 在 $k$ 上\Term{separable-extension}{可分}，则 $F(k)\to F(A)$ 为\Term{injective}{单射}。由对 $k$ 和 $K$ 所作的假设，存在一个 $k$ \Term{homomorphism}{同态} $A\to k$；于是复合 $F(k)\to F(A)\to F(k)$ 为\Term{identity-map}{恒等映射}，结论随即成立。在 b) 中，同样经取\Term{limit}{极限}，归结到证明：若 $A$ 是 $k$ 上的\Term{finite-type}{有限型}\Term{integral-algebra}{整代数}，其\Term{fraction-field}{分式域}是 $k$ 的\Term{pure-transcendental-extension}{纯超越扩张}，则 $F(k)\to F(A)$ 为\Term{injective}{单射}。由于 $k$ 无限，此假设仍蕴含存在一个 $k$ \Term{homomorphism}{同态} $A\to k$，于是如前完成证明。
\textbf{7.18.3.} 现在假设 $F$ 是取值于\Term{abelian-group}{阿贝尔群}\Term{category}{范畴} $(\mathrm{Ab})$ 的\Term{functor}{函子}，并假设对 $k$ 的每个\Term{finite-extension}{有限扩张} $k'$，定义了\Term{norm-homomorphism}{范数同态}
\[N_{k'/k}:F(k')\longrightarrow F(k),\]
使得复合
\[F(k)\longrightarrow F(k')\xrightarrow{N_{k'/k}}F(k)\]
是乘以 $[k':k]$。在这些条件下，有：
\SourcePage{602}
\begin{enumerate}
\item[(i)] 对 $k$ 的任意\Term{extension}{扩张} $K$，$F(k)\to F(K)$ 的核为\Term{torsion-group}{挠群}。
\item[(ii)] 对 $k$ 的任意\Term{pure-transcendental-extension}{纯超越扩张} $K$，\Term{homomorphism}{同态} $F(k)\to F(K)$ 为\Term{injective}{单射}。（注意，这里不必假设 $k$ 无限。）
\end{enumerate}
在 (i) 中，只须证明 $F(k)\to F(A)$ 的核中的每个元素 $x$ 都是\Term{torsion}{挠元}，其中 $A$ 是 $k$ 上的\Term{finite-type}{有限型}\Term{integral-algebra}{整代数}。利用一个 $k$ \Term{homomorphism}{同态} $A\to k'$，其中 $k'$ 为\Term{finite}{有限} $k$ \Term{algebra}{代数}，可归结到 $A=k'$ 是 $k$ 的\Term{finite-extension}{有限扩张}的情形。此时由前面对\Term{norm-homomorphism}{范数同态}要求的性质，有 $[k':k]x=0$，这证明 $x$ 是\Term{torsion}{挠元}。如果作为 $A$ 的商的各\Term{finite-extension}{有限扩张} $k'/k$ 的\Term{degree}{次数}的最大公因数为 $1$，则甚至有 $x=0$。容易看出，若 $A$ 的\Term{fraction-field}{分式域} $K$ 是 $k$ 的\Term{pure-transcendental-extension}{纯超越扩张}，这一条件就成立（只有 $k$ \Term{finite}{有限}时需要另作验证）。由此即得 (ii) 的结论。
\textbf{7.18.4.} 假设对于\Term{affine-space}{仿射空间} $\operatorname{Spec}k[T_1,\ldots,T_n]$ 的每个\Term{nonempty}{非空}\Term{affine-open}{仿射开集} $U$，若 $A$ 是 $U$ 的\Term{ring}{环}，则\Term{homomorphism}{同态} $F(k)\to F(A)$ 为\Term{surjective}{满射}（比较第 IX 讲 1.6\SGAsourcegap{原页“IX 1.6”后未填完整引注编号}）。于是利用 $F$ 与\Term{filtered-inductive-limit}{滤过归纳极限}交换，得到：对于 $k$ 的每个\Term{pure-transcendental-extension}{纯超越扩张} $K$，$F(k)\to F(K)$ 为\Term{surjective}{满射}。若还假设 $k$ 无限，或者能像 7.18.3 那样定义\Term{norm-homomorphism}{范数同态}，则得出这一\Term{homomorphism}{同态}甚至为\Term{bijection}{双射}。
\textbf{7.18.5.} 举例来说，考虑 $k$ 上的\Term{finite-type}{有限型}\Term{scheme}{概形} $X$，令
\[F(A)=K_\bullet(X\otimes_k A).\]
\SourcePage{603}
由第 IV 讲 2.12，这是\Term{category}{范畴} $C$ 上的一个\Term{functor}{函子}。由第 IV 讲 3.2.3，它与\Term{filtered-inductive-limit}{滤过归纳极限}交换。此外，若 $k'$ 是 $k$ 的\Term{finite-extension}{有限扩张}，则\Term{projection}{投影} $p:X\otimes_k k'\to X$ 是\Term{finite-morphism}{有限}、\Term{flat-morphism}{平坦}、\Term{finite-presentation}{有限表现}的\Term{morphism}{态射}，因而定义\Term{direct-image}{正像}\Term{homomorphism}{同态}
\[p_*=N_{k'/k}:K_\bullet(X\otimes_k k')\longrightarrow K_\bullet(X).\]
立即可验证，7.18.3 中对\Term{norm-homomorphism}{范数同态}要求的条件确实满足。最后，若 $A$ 是\Term{affine-space}{仿射空间} $\operatorname{Spec}k[T_1,\ldots,T_n]$ 的一个\Term{open-subscheme}{开集} $U$ 的\Term{affine-algebra}{仿射代数}，则 $F(k)\to F(A)$ 等同于映射 $K_\bullet(X)\to K_\bullet(X\times_k U)$；当 $X$ \Term{noetherian}{诺特}时，由第 IX 讲 1.6\SGAsourcegap{原页“IX 1.6”后未填完整引注编号}，它是\Term{surjective}{满射}。利用以上证明的结果，得到：
\begin{enumerate}
\item[(i)] 若 $K/k$ 是 $k$ 的\Term{pure-transcendental-extension}{纯超越扩张}，则 $K_\bullet(X)\to K_\bullet(X_K)$ 为\Term{injective}{单射}；若 $X$ \Term{noetherian}{诺特}，它甚至为\Term{bijection}{双射}。
\item[(ii)] 若 $k$ \Term{separably-closed}{可分闭}，且 $K/k$ 为\Term{separable-extension}{可分扩张}，则 $K_\bullet(X)\to K_\bullet(X_K)$ 为\Term{injective}{单射}。
\item[(iii)] 对 $k$ 的任意\Term{extension}{扩张} $K$，$K_\bullet(X)\to K_\bullet(X_K)$ 的核为\Term{torsion-group}{挠群}。
\end{enumerate}
如果仅假设 $X$ \Term{noetherian}{诺特}，只要限于 $k$ 的\Term{finite-type}{有限型}\Term{extension}{扩张} $K$（从而 $X_K$ 仍为\Term{finite-type}{有限型}），这些结果仍成立。此外，将\Term{functor}{函子} $K_\bullet$ 换成 $K'$，实质上得到同样的结果，唯一区别是：在 (i) 中，\Term{generic}{一般}不能断言为\Term{bijection}{双射}，因为没有将 $K_\bullet$ 换成 $K'$ 的\Term{homotopy-theorem}{同伦定理}；为得出\Term{bijection}{双射}性，必须假设 $X$ \Term{noetherian}{诺特}\Term{regular-scheme}{正则}，这就意味着对 $X$ 及 $X_K$，\Term{group}{群} $K'$、$K_\bullet$ 相同。
\SourcePage{604}
将 $K_\bullet$ 或 $K'$ 换成 $\operatorname{Gr}'$，实质上仍得到同样的结果，但对于 (i) 中的\Term{bijection}{双射}性断言，有同样的保留。如果仅考虑 $\operatorname{Gr}^1=\operatorname{Pic}$，则当 $X$ \Term{noetherian}{诺特}且\Term{normal-scheme}{正规}时，此\Term{bijection}{双射}性断言成立，因为此时对 $k$ 上\Term{affine-space}{仿射空间}的任一\Term{open-subscheme}{开集} $U$，立即可知 $\operatorname{Pic}(X)\to\operatorname{Pic}(X\times_k U)$ 为\Term{surjective}{满射}（EGA $\operatorname{Err}_{\mathrm{IV}}$ 53，21.4.13）。
\textbf{参考文献}
\begin{enumerate}
\item[{[1]}] E. Dickson，\emph{\Term{linear}{线性}\Term{group}{群}}（\emph{Linear groups}），Dover，New York，1958。
\item[{[2]}] A. Grothendieck，TDTE，V、VI；Séminaire Bourbaki，232、236。
\item[{[3]}] S. Kleiman，\emph{迈向\Term{ample}{丰沛}性的数值理论}（\emph{Towards a numerical theory of ampleness}），Ann. of Math. \textbf{84} (1966)，293--344 页。
\item[{[4]}] J. Murre，\emph{从域上的\Term{prescheme}{预概形}\Term{category}{范畴}到阿贝尔\Term{group}{群}\Term{category}{范畴}的\Term{contravariant-functor}{反变函子}}（\emph{On contravariant functors from the category of préschémes over a field into the category of abelian groups}），Publications Mathématiques de l'IHES，23，1964。
\item[{[5]}] C. Chevalley，\emph{\Term{rational-equivalence}{有理等价}类}（\emph{Les classes d'équivalence rationnelles}），I、II，Séminaire à l'ENS 1955，第 2、3 讲。
\item[{[6]}] S. Abhyankar，\emph{算术\Term{surface}{曲面}的\Term{resolution-singularities}{奇点消解}}（\emph{Resolution of singularities of arithmetical surfaces}），载于 \emph{Arithmetical Algebraic Geometry}，Purdue University，Harper \& Row，New York，1965。
\end{enumerate}
